\section{Resultados e Discussão}
\label{sec:experimentos}

Esta seção detalha os experimentos empíricos conduzidos com a arquitetura proposta, avaliando convergência de função de perda, taxa de acerto (\textit{accuracy}) e estabilidade numérica.

\subsection{Ambiente Experimental}
Os experimentos foram executados utilizando o ambiente integrado em Python, com dados normalizados e divisão entre conjuntos de treino e validação.

\begin{table}[ht]
\centering
\caption{Resumo dos hiperparâmetros utilizados nos ensaios experimentais.}
\label{tab:hiperparametros}
\begin{tabular}{@{}ll@{}}
\toprule
\textbf{Parâmetro} & \textbf{Valor Adotado} \\ \midrule
Taxa de Aprendizado ($\eta$) & 0.01 -- 0.1 \\
Otimizador & Descida do Gradiente Estocástico (SGD) \\
Função de Custo & Erro Quadrático Médio (MSE) / Entropia Cruzada \\
Épocas de Treinamento & 500 \\ \bottomrule
\end{tabular}
\end{table}

Os resultados preliminares evidenciam a convergência monótona da perda e a consistência dos gradientes calculados através da estrutura em grafo.
